% =====================================================================
% Notas técnicas: Apartado 4 (CAP y PACELC) y Apartado 7 (Postura crítica)
% Responsable y expositor: Mendoza Villarroel Elian Jossue
% Se incluye en notas/notas.tex con % =====================================================================
% Notas técnicas: Apartado 4 (CAP y PACELC) y Apartado 7 (Postura crítica)
% Responsable y expositor: Mendoza Villarroel Elian Jossue
% Se incluye en notas/notas.tex con % =====================================================================
% Notas técnicas: Apartado 4 (CAP y PACELC) y Apartado 7 (Postura crítica)
% Responsable y expositor: Mendoza Villarroel Elian Jossue
% Se incluye en notas/notas.tex con % =====================================================================
% Notas técnicas: Apartado 4 (CAP y PACELC) y Apartado 7 (Postura crítica)
% Responsable y expositor: Mendoza Villarroel Elian Jossue
% Se incluye en notas/notas.tex con \input{notas/apartados/cap_comparacion}
% Claves del .bib: ahmed2023consistency, lee2023consistency,
%                  gilbert2012perspectives, lima2023causal
% Versión inicial: falta el sistema analizado, la comparación y el aporte al PFC.
% =====================================================================

\subsection*{Notas técnicas: CAP y PACELC}

\paragraph{Teorema CAP.}
El teorema CAP establece un compromiso fundamental entre tres propiedades de
los sistemas distribuidos replicados: consistencia, disponibilidad y
tolerancia a particiones \cite[p.~172]{ahmed2023consistency}. Cuando ocurre
una interrupción en la comunicación entre sus nodos, no es posible garantizar
simultáneamente la consistencia y la disponibilidad. Por esta razón, si el
sistema debe seguir funcionando a pesar de una partición de red, tendrá que
priorizar una de estas dos propiedades \cite[p.~2]{lee2023consistency}.

\paragraph{Consistencia (C).}
En el sentido de CAP, la consistencia fuerte exige que las operaciones se
observen como si ocurrieran en una sola copia y en un orden compatible con
el tiempo real \cite[Sec.~2]{gilbert2012perspectives}. Se trata de una
propiedad de seguridad: las respuestas deben respetar esa garantía
\cite[Sec.~1]{gilbert2012perspectives}.

\paragraph{Disponibilidad (A) y tolerancia a particiones (P).}
La disponibilidad es la capacidad de responder a las solicitudes de los
usuarios que leen o modifican el estado compartido
\cite[p.~2]{lee2023consistency}. La tolerancia a particiones parte de que la
comunicación entre servidores no es fiable: los mensajes pueden retrasarse o
perderse y los servidores pueden quedar divididos en grupos que no se
comunican \cite[Sec.~2]{gilbert2012perspectives}.

\paragraph{Comportamiento durante una partición.}
Cuando una réplica rechaza una solicitud o no responde dentro del tiempo de
espera, la base de datos debe decidir entre devolver un resultado
posiblemente obsoleto, sacrificando la consistencia, o informar que el
servicio no está disponible. Los sistemas CP, como HBase o BigTable,
renuncian a la disponibilidad; los sistemas AP, como Cassandra, Dynamo o
Riak, relajan la consistencia \cite[p.~172]{ahmed2023consistency}.

\paragraph{Teorema PACELC.}
PACELC plantea dos decisiones en sistemas replicados. Si hay una partición,
el compromiso es entre disponibilidad y consistencia. Si no hay una
partición, el compromiso es entre latencia y consistencia
\cite{ahmed2023consistency}. Un caso del segundo compromiso es el acceso
causal en almacenamiento georreplicado, donde esperar a que se sincronicen
las réplicas puede retrasar el acceso \cite{lima2023causal}.

\subsection*{Postura crítica}

Priorizar consistencia o disponibilidad depende de lo que necesite cada
aplicación. En un cruce controlado por vehículos autónomos conviene la
consistencia, porque todos deben compartir la misma información antes de
avanzar; en cambio, en el sistema de frenado conviene la disponibilidad,
porque la reacción tiene que ser inmediata
\cite[pp.~2--3]{lee2023consistency}. En una tienda en línea, una consulta
puede devolver datos desactualizados, mientras que la compra requiere
consistencia para que el pedido refleje lo que eligió el cliente
\cite[Sec.~4.4]{gilbert2012perspectives}. En consecuencia, ningún sistema es
superior en todos los casos: cada uno equilibra estas propiedades según su
propósito.

% Pendiente para próximos commits:
%  - Sistema analizado: qué prioriza con y sin partición, clasificación
%    PACELC y ejemplo técnico de su documentación oficial.
%  - Comparación con un segundo sistema real (tabla del apartado 7).
%  - Aporte al PFC para el apartado 8.

% Claves del .bib: ahmed2023consistency, lee2023consistency,
%                  gilbert2012perspectives, lima2023causal
% Versión inicial: falta el sistema analizado, la comparación y el aporte al PFC.
% =====================================================================

\subsection*{Notas técnicas: CAP y PACELC}

\paragraph{Teorema CAP.}
El teorema CAP establece un compromiso fundamental entre tres propiedades de
los sistemas distribuidos replicados: consistencia, disponibilidad y
tolerancia a particiones \cite[p.~172]{ahmed2023consistency}. Cuando ocurre
una interrupción en la comunicación entre sus nodos, no es posible garantizar
simultáneamente la consistencia y la disponibilidad. Por esta razón, si el
sistema debe seguir funcionando a pesar de una partición de red, tendrá que
priorizar una de estas dos propiedades \cite[p.~2]{lee2023consistency}.

\paragraph{Consistencia (C).}
En el sentido de CAP, la consistencia fuerte exige que las operaciones se
observen como si ocurrieran en una sola copia y en un orden compatible con
el tiempo real \cite[Sec.~2]{gilbert2012perspectives}. Se trata de una
propiedad de seguridad: las respuestas deben respetar esa garantía
\cite[Sec.~1]{gilbert2012perspectives}.

\paragraph{Disponibilidad (A) y tolerancia a particiones (P).}
La disponibilidad es la capacidad de responder a las solicitudes de los
usuarios que leen o modifican el estado compartido
\cite[p.~2]{lee2023consistency}. La tolerancia a particiones parte de que la
comunicación entre servidores no es fiable: los mensajes pueden retrasarse o
perderse y los servidores pueden quedar divididos en grupos que no se
comunican \cite[Sec.~2]{gilbert2012perspectives}.

\paragraph{Comportamiento durante una partición.}
Cuando una réplica rechaza una solicitud o no responde dentro del tiempo de
espera, la base de datos debe decidir entre devolver un resultado
posiblemente obsoleto, sacrificando la consistencia, o informar que el
servicio no está disponible. Los sistemas CP, como HBase o BigTable,
renuncian a la disponibilidad; los sistemas AP, como Cassandra, Dynamo o
Riak, relajan la consistencia \cite[p.~172]{ahmed2023consistency}.

\paragraph{Teorema PACELC.}
PACELC plantea dos decisiones en sistemas replicados. Si hay una partición,
el compromiso es entre disponibilidad y consistencia. Si no hay una
partición, el compromiso es entre latencia y consistencia
\cite{ahmed2023consistency}. Un caso del segundo compromiso es el acceso
causal en almacenamiento georreplicado, donde esperar a que se sincronicen
las réplicas puede retrasar el acceso \cite{lima2023causal}.

\subsection*{Postura crítica}

Priorizar consistencia o disponibilidad depende de lo que necesite cada
aplicación. En un cruce controlado por vehículos autónomos conviene la
consistencia, porque todos deben compartir la misma información antes de
avanzar; en cambio, en el sistema de frenado conviene la disponibilidad,
porque la reacción tiene que ser inmediata
\cite[pp.~2--3]{lee2023consistency}. En una tienda en línea, una consulta
puede devolver datos desactualizados, mientras que la compra requiere
consistencia para que el pedido refleje lo que eligió el cliente
\cite[Sec.~4.4]{gilbert2012perspectives}. En consecuencia, ningún sistema es
superior en todos los casos: cada uno equilibra estas propiedades según su
propósito.

% Pendiente para próximos commits:
%  - Sistema analizado: qué prioriza con y sin partición, clasificación
%    PACELC y ejemplo técnico de su documentación oficial.
%  - Comparación con un segundo sistema real (tabla del apartado 7).
%  - Aporte al PFC para el apartado 8.

% Claves del .bib: ahmed2023consistency, lee2023consistency,
%                  gilbert2012perspectives, lima2023causal
% Versión inicial: falta el sistema analizado, la comparación y el aporte al PFC.
% =====================================================================

\subsection*{Notas técnicas: CAP y PACELC}

\paragraph{Teorema CAP.}
El teorema CAP establece un compromiso fundamental entre tres propiedades de
los sistemas distribuidos replicados: consistencia, disponibilidad y
tolerancia a particiones \cite[p.~172]{ahmed2023consistency}. Cuando ocurre
una interrupción en la comunicación entre sus nodos, no es posible garantizar
simultáneamente la consistencia y la disponibilidad. Por esta razón, si el
sistema debe seguir funcionando a pesar de una partición de red, tendrá que
priorizar una de estas dos propiedades \cite[p.~2]{lee2023consistency}.

\paragraph{Consistencia (C).}
En el sentido de CAP, la consistencia fuerte exige que las operaciones se
observen como si ocurrieran en una sola copia y en un orden compatible con
el tiempo real \cite[Sec.~2]{gilbert2012perspectives}. Se trata de una
propiedad de seguridad: las respuestas deben respetar esa garantía
\cite[Sec.~1]{gilbert2012perspectives}.

\paragraph{Disponibilidad (A) y tolerancia a particiones (P).}
La disponibilidad es la capacidad de responder a las solicitudes de los
usuarios que leen o modifican el estado compartido
\cite[p.~2]{lee2023consistency}. La tolerancia a particiones parte de que la
comunicación entre servidores no es fiable: los mensajes pueden retrasarse o
perderse y los servidores pueden quedar divididos en grupos que no se
comunican \cite[Sec.~2]{gilbert2012perspectives}.

\paragraph{Comportamiento durante una partición.}
Cuando una réplica rechaza una solicitud o no responde dentro del tiempo de
espera, la base de datos debe decidir entre devolver un resultado
posiblemente obsoleto, sacrificando la consistencia, o informar que el
servicio no está disponible. Los sistemas CP, como HBase o BigTable,
renuncian a la disponibilidad; los sistemas AP, como Cassandra, Dynamo o
Riak, relajan la consistencia \cite[p.~172]{ahmed2023consistency}.

\paragraph{Teorema PACELC.}
PACELC plantea dos decisiones en sistemas replicados. Si hay una partición,
el compromiso es entre disponibilidad y consistencia. Si no hay una
partición, el compromiso es entre latencia y consistencia
\cite{ahmed2023consistency}. Un caso del segundo compromiso es el acceso
causal en almacenamiento georreplicado, donde esperar a que se sincronicen
las réplicas puede retrasar el acceso \cite{lima2023causal}.

\subsection*{Postura crítica}

Priorizar consistencia o disponibilidad depende de lo que necesite cada
aplicación. En un cruce controlado por vehículos autónomos conviene la
consistencia, porque todos deben compartir la misma información antes de
avanzar; en cambio, en el sistema de frenado conviene la disponibilidad,
porque la reacción tiene que ser inmediata
\cite[pp.~2--3]{lee2023consistency}. En una tienda en línea, una consulta
puede devolver datos desactualizados, mientras que la compra requiere
consistencia para que el pedido refleje lo que eligió el cliente
\cite[Sec.~4.4]{gilbert2012perspectives}. En consecuencia, ningún sistema es
superior en todos los casos: cada uno equilibra estas propiedades según su
propósito.

% Pendiente para próximos commits:
%  - Sistema analizado: qué prioriza con y sin partición, clasificación
%    PACELC y ejemplo técnico de su documentación oficial.
%  - Comparación con un segundo sistema real (tabla del apartado 7).
%  - Aporte al PFC para el apartado 8.

% Claves del .bib: ahmed2023consistency, lee2023consistency,
%                  gilbert2012perspectives, lima2023causal
% Versión inicial: falta el sistema analizado, la comparación y el aporte al PFC.
% =====================================================================

\subsection*{Notas técnicas: CAP y PACELC}

\paragraph{Teorema CAP.}
El teorema CAP establece un compromiso fundamental entre tres propiedades de
los sistemas distribuidos replicados: consistencia, disponibilidad y
tolerancia a particiones \cite[p.~172]{ahmed2023consistency}. Cuando ocurre
una interrupción en la comunicación entre sus nodos, no es posible garantizar
simultáneamente la consistencia y la disponibilidad. Por esta razón, si el
sistema debe seguir funcionando a pesar de una partición de red, tendrá que
priorizar una de estas dos propiedades \cite[p.~2]{lee2023consistency}.

\paragraph{Consistencia (C).}
En el sentido de CAP, la consistencia fuerte exige que las operaciones se
observen como si ocurrieran en una sola copia y en un orden compatible con
el tiempo real \cite[Sec.~2]{gilbert2012perspectives}. Se trata de una
propiedad de seguridad: las respuestas deben respetar esa garantía
\cite[Sec.~1]{gilbert2012perspectives}.

\paragraph{Disponibilidad (A) y tolerancia a particiones (P).}
La disponibilidad es la capacidad de responder a las solicitudes de los
usuarios que leen o modifican el estado compartido
\cite[p.~2]{lee2023consistency}. La tolerancia a particiones parte de que la
comunicación entre servidores no es fiable: los mensajes pueden retrasarse o
perderse y los servidores pueden quedar divididos en grupos que no se
comunican \cite[Sec.~2]{gilbert2012perspectives}.

\paragraph{Comportamiento durante una partición.}
Cuando una réplica rechaza una solicitud o no responde dentro del tiempo de
espera, la base de datos debe decidir entre devolver un resultado
posiblemente obsoleto, sacrificando la consistencia, o informar que el
servicio no está disponible. Los sistemas CP, como HBase o BigTable,
renuncian a la disponibilidad; los sistemas AP, como Cassandra, Dynamo o
Riak, relajan la consistencia \cite[p.~172]{ahmed2023consistency}.

\paragraph{Teorema PACELC.}
PACELC plantea dos decisiones en sistemas replicados. Si hay una partición,
el compromiso es entre disponibilidad y consistencia. Si no hay una
partición, el compromiso es entre latencia y consistencia
\cite{ahmed2023consistency}. Un caso del segundo compromiso es el acceso
causal en almacenamiento georreplicado, donde esperar a que se sincronicen
las réplicas puede retrasar el acceso \cite{lima2023causal}.

\subsection*{Postura crítica}

Priorizar consistencia o disponibilidad depende de lo que necesite cada
aplicación. En un cruce controlado por vehículos autónomos conviene la
consistencia, porque todos deben compartir la misma información antes de
avanzar; en cambio, en el sistema de frenado conviene la disponibilidad,
porque la reacción tiene que ser inmediata
\cite[pp.~2--3]{lee2023consistency}. En una tienda en línea, una consulta
puede devolver datos desactualizados, mientras que la compra requiere
consistencia para que el pedido refleje lo que eligió el cliente
\cite[Sec.~4.4]{gilbert2012perspectives}. En consecuencia, ningún sistema es
superior en todos los casos: cada uno equilibra estas propiedades según su
propósito.

% Pendiente para próximos commits:
%  - Sistema analizado: qué prioriza con y sin partición, clasificación
%    PACELC y ejemplo técnico de su documentación oficial.
%  - Comparación con un segundo sistema real (tabla del apartado 7).
%  - Aporte al PFC para el apartado 8.
